\documentclass[12pt]{article}
\usepackage[T2A]{fontenc}
\usepackage[utf8]{inputenc}
\usepackage[russian]{babel}
\usepackage[a4paper,left=2.5cm,right=2.5cm,top=2cm,bottom=2cm]{geometry}
\usepackage{indentfirst}
\usepackage{setspace}

\setlength{\parindent}{1.25em}
\onehalfspacing

% Запрет некрасивых переносов («филь-трах», «физи-ка» и т.п.)
\hyphenation{физика математика фильтр фильтра фильтре фильтрах фильтров
  фильтрами фильтрации}

\begin{document}

\begin{center}
  {\large\bfseries ОТЗЫВ НАУЧНОГО РУКОВОДИТЕЛЯ}\\[2pt]
  {\large\bfseries на выпускную квалификационную работу}
\end{center}

\medskip

\noindent Студент: Сажнев Максим Евгеньевич\\
Тема ВКР: «Использование аппроксимационных умножителей в цифровых
фильтрах»\\
Направление подготовки / специальность: 03.03.01 Прикладные математика и
физика\\
Направленность (профиль) подготовки: Радиотехника и компьютерные технологии

\medskip

В условиях роста требований к энергоэффективности вычислительных систем,
встраиваемых устройств и специализированных аппаратных ускорителей
исследование методов оптимизации арифметических блоков представляет
существенный научный и практический интерес. В связи с этим выпускная
квалификационная работа, посвящённая снижению энергопотребления цифровых
умножителей, применяемых в адаптивных КИХ-фильтрах, является актуальной.

В работе предложен и исследован подход к снижению энергопотребления
умножителя Бута по основанию~4 за счёт вынесения блока кодировки за пределы
основного умножителя и применения алгоритма перекодировки, основанного на
избыточности кода Бута. При этом была учтена особенность применения данного
подхода в задачах адаптивной цифровой фильтрации, когда коэффициенты
обновляются существенно реже, чем входные данные. Использование такого
подхода позволило снизить энергопотребление умножителя примерно на 4\% при
сохранении приемлемой точности вычислений без увеличения площади умножителя
при аппаратной реализации.

Достоверность результатов подтверждается численным моделированием, а также
экспериментальными результатами синтеза RTL-реализации предложенной
архитектуры и базовой реализации умножителя Бута на языке Verilog с
использованием академической библиотеки стандартных ячеек.

Теоретическая значимость работы состоит в адаптации метода перекодировки
частичных произведений к архитектуре параллельного умножителя, применяемого
в составе адаптивного фильтра. Практическая значимость заключается в
предложении аппаратной архитектуры, которая может быть использована при
проектировании энергоэффективных цифровых фильтров и специализированных
вычислительных блоков.

К положительным сторонам работы следует отнести достаточный объём
экспериментального материала, включающий оценку площади, мощности, точности
и масштабируемости по разрядности, полученный в результате аппаратного
синтеза RTL-реализации предложенной архитектуры умножителя.

К недостаткам работы можно отнести то, что экспериментальная оценка
результатов синтеза проводилась в рамках реализации адаптивного КИХ-фильтра,
а не всей системы адаптивной коррекции модуля цифровой обработки сигналов.
Указанные замечания не снижают общей высокой оценки выполненной работы.

В ходе выполнения выпускной квалификационной работы студент М.Е.~Сажнев
проявил себя как самостоятельный и подготовленный исследователь. Он
продемонстрировал хорошие навыки работы с литературой, владение методами
цифрового проектирования, способность проводить экспериментальные
исследования и анализировать полученные результаты. Объём проанализированного
материала и качество выполненной реализации свидетельствуют о
сформированности необходимых профессиональных компетенций. Сажнев М.Е. также
показал умение ясно представлять результаты исследования и аргументировать
сделанные выводы.

Считаю, что выпускная квалификационная работа Сажнева М.Е. соответствует
требованиям, установленным к ВКР студентов МФТИ, заслуживает оценки
«отлично», а её автор Сажнев Максим Евгеньевич заслуживает присвоения
квалификации бакалавра по соответствующему направлению подготовки.

\bigskip

\noindent Научный руководитель:

\medskip

\noindent к.т.н., доцент\hfill Бахурин С.~А.\hspace{2em}

\smallskip

\noindent\hspace{\dimexpr0.38\textwidth+1cm\relax}(подпись)

\bigskip\bigskip

\noindent «\hspace{2em}»\hspace{4em}2026 г.

\end{document}
